The second implementation stage addresses the first release's engineering
gaps. Its source and benchmark records are separate from the historical
snapshots in Section~\ref{sec:computation}. The mathematical distinction
between validity, termination and practical cost remains essential:
an accepted rational certificate is exact, whereas a resource-limited
search can finish without finding one.

\subsection{Shared finite tables and automatic decompositions}
The reusable finite-tree engine computes both message directions, all
coordinate min-marginals and attaining assignments for finite factor tables,
including infeasible entries. It checks tree structure and running
intersection; the model adapters check factor coverage. Its callers build
different objective and feasibility
tables but share this computational rule. Certificate checkers separately
validate the model-specific mathematical bounds. The base and constrained
checkers independently implement Bellman replay; the convex-value checker
reuses the finite-tree engine after independently checking its local
optimization witnesses.

When a box-QP decomposition is omitted, a deterministic minimum-fill
heuristic constructs one from the interaction graph. A supplied
decomposition remains supported and is checked. The resolved bags and tree
are serialized, so certificate replay does not depend on reproducing the
heuristic. Returned width is an upper bound, not a minimum-treewidth claim.
Table budgets are checked before oversized allocations. Sparse neighbor
lists reduce zero-product work in coordinate polishing and gradients;
dense model storage and exact rational arithmetic remain practical costs.

\subsection{Exact rational output in the general grid solver}
The \texttt{solve\_exact} wrapper repeatedly requests finer original-model
grid bounds, generates rational candidates, and applies
Proposition~\ref{prop:candidateheight}. The height bound is computed once
from the original coefficients and bounds. The implementation sharpens
\eqref{eq:height} using a product of integer Hessian row norms, which bounds
every relevant principal determinant by Hadamard's inequality. Neither
filtered hull endpoints nor candidate denominators alter that original
optimum-height bound.

Candidates include the current feasible incumbent, bounded-denominator
coordinate reconstruction, and exact stationarity on a proposed active
face. Singular stationary systems use rational linear feasibility.
These are candidate-generation procedures: stationarity or proximity
alone never establishes global optimality. The separate exact-output
checker replays the underlying original-domain certificate, verifies the
candidate against the original model, recomputes its objective and reduced
denominator, and checks the strict rational-separation inequality.

The wrapper exposes round, stage, table and time limits, with configurable
initial precision.
An interrupted run retains its certified enclosure and says that exact
output was not reached. Theorem~\ref{thm:generalfinite} now proves eventual exact output for
arbitrary nonunique rational box QPs with sufficiently expanded resource
caps. The stronger point-growth running-time theorem retains the
uniqueness hypothesis stated in Section~\ref{sec:exact}. General finite
termination does not provide that stronger efficiency guarantee.

\subsection{Reusable exact linear and convex quadratic optimization}
The shared rational optimization module implements two-phase simplex with
Bland's pivot rule. Successful LP results carry primal/dual witnesses,
infeasibility carries a Farkas witness, and unboundedness carries a feasible
point and improving ray. The result checker verifies the original rational
constraints and algebraic witness without replaying simplex pivots.
A pivot limit is inconclusive; it is not an infeasibility certificate.

The convex box-QP oracle checks positive semidefiniteness exactly and
searches active faces with exact linear algebra and LP feasibility for
singular faces. A feasible point with checked KKT signs certifies the global
minimum of the convex quadratic. Face and pivot limits bound this
reference implementation. Neither simplex nor active-face enumeration is
claimed here as a polynomial-time implementation of the theoretical
convex-QP oracle. Their purpose is to make the restricted algorithms
executable while retaining independently checkable answers.

\subsection{Original-model recourse composition}
The recourse pipeline recognizes globally affine convex responses,
substitutes accepted responses, solves the reduced model, and lifts the
result to the original coordinates. Each accepted response has a checked
positive-semidefinite private Hessian, feasible whole-box affine range,
and whole-box KKT identities and signs. Thus the original and reduced
optimum values coincide. The final certificate stores every reduction,
including fixed-coordinate substitutions, so a checker can reconstruct
the chain and verify the final incumbent in the original model.

Candidate private blocks can be supplied or proposed from continuous
components, private leaf-bag variables and single coordinates. This
selection is heuristic. A rejected block or exhausted recognition budget
remains in the model, and a failure to discover a reduction does not prove
that no useful partition exists. The reduced problem can use the sparse
grid solver or signed-concave minimum-cut recourse. The latter adapter
checks signs and separates a small core before applying the existing exact
conditional-cut search. The full reduction and search cost belongs in the
reported total; preprocessing is not a free input promise.

Changing-active-set convex recourse is also executable. Supplied disjoint
continuous PSD blocks, or heuristically discovered blocks, become local
conditional value factors. The implementation caches rational convex-QP
values and KKT witnesses at attachment grid points, uses corrected geometric
grids on the retained coordinates, and supports exact original-QP height
acceptance. Private feasible sets in this implementation are boxes; the
general private-polytope theorem has a wider input class. The checker
verifies local PSD/KKT data, recomputes finite-tree values using the shared
engine, and checks the original-domain history. It invokes no LP or QP
optimizer.

For scalar attachments, the integrated optional constructor enumerates
private active patterns within a stated budget, certifies a complete
interval partition, and uses its maximal private reduced curvature to
sharpen the retained diagonal. Contributions from multiple scalar blocks
add without a global overlay. The current constructor requires a positive-
definite private Hessian and nonfixed intervals; unsupported or over-budget
blocks retain the safe direct bound. The default pattern cap is $81$.
In tested changing-response examples this changes a retained bound from
$2M+6$ to $6$ for stiffness $M=1$ and $M=100$. General multidimensional
short-partition construction remains outside the implementation.

The mixed convex/concave submodular backend goes beyond minimum cuts.
After verified sign changes, endpoint reduction of separately concave
coordinates leaves a submodular set function evaluated by conditional
convex QPs. Rational LP cutting planes minimize its Lov\'asz extension;
greedy chains supply separating affine functions. The final lower bound
uses a finite convex mixture of checked greedy base vectors, and a feasible
conditional minimizer attains it. Its independent checker verifies the
conditional KKT witnesses and the mixture bound. This is an exact capped
reference implementation, without a polynomial iteration guarantee for
its cutting-plane search. It stores completed cached query witnesses, so
its certificate storage follows the query budget; the implementation does
not claim the theorem's polynomial certificate-size bound.


\subsection{Coupled constraints and optimal sets}
The reusable TU solver uses feasible common-mesh tables on general tree
decompositions, all coordinate min-marginals, and complete filtering
history. Constraint scopes join objective scopes in the decomposition.
The input includes an exactly checked TU certificate: supported structural
forms include signed incidence and consecutive-ones matrices; a bounded
all-minors fallback handles small matrices. A refused TU verification is
not a proof of non-TU structure. Original integral alignment, finite
integer labels and verified continuous curvature remain assumptions.
The optional union representation preserves separated surviving cells.
Exact output combines original-face rational candidates, the generic
stationary-face recovery LP and a height-separation check. With resource
caps removed, the recovery construction also covers arbitrary nonunique
quadratic minima; only the finite-projection subclass has the stronger
state bound. The width-dependent exponent and initial capacity costs
stated in Section~\ref{sec:extensions} are unchanged.

The optimal-set module constructs a compact Bellman-residual descriptor
for the entire optimal set of a coordinatewise concave mixed-box QP. A
separate checker verifies the recurrence and the support conditions,
including possible interior integer labels. For the diagonal-certifiable
continuous class, the implementation actually runs finite proximal trials
with guessed growth ratios, performs rational stationary recovery, and
accepts only the original KKT and positive-semidefinite certificate. A
false growth guess cannot justify acceptance. Exhausting the configured
search limits yields an inconclusive result, not a negative membership
decision for the class.

\subsection{Explicit polynomials and exact implicit boundary output}
The polynomial implementation accepts explicit rational monomial factors,
constructs or verifies their sparse decomposition, and computes upper
coordinate-curvature bounds by exact interval arithmetic. It uses the shared
finite-tree engine, corrected grids and complete pruning histories. Native
integer output has the coefficient-denominator lattice stopping rule from
Section~\ref{sec:polynomial}, computed after substituting any fixed rational
coordinates. Its checker independently recomputes that denominator and the
strict pre-acceptance gap. Continuous approximation does not round an
algebraic optimum to a rational exact answer.

The boundary search feeds finer globally certified polynomial boxes to a
separate exact acceptance test. It first fixes native labels, then
repeatedly certifies derivative signs and clamps the corresponding original
bounds. Such a clamp preserves at least one global optimizer; a weak sign
need not preserve every optimizer. On the remaining face, let $H_0$ be the
midpoint Hessian and let $\delta$ bound the maximum absolute row sum of
$H(x)-H_0$ throughout the box. An exact positive-definiteness check on
$H_0-\delta I$ proves uniform strong convexity. The polynomial together
with this rational face box is then an exact finite descriptor of one
global optimizer, possibly with irrational coordinates and value.

The boundary verifier independently recomputes monomial derivative
intervals, reductions and the matrix bound, after replaying the global
grid proof. It does not optimize the patch. The tested polynomial
$(x^2-1/2)^2+y+xy^2$ yields a certified patch for
$(x,y)=(1/\sqrt2,0)$ after four accuracy rounds. A bounded restart search
implements discovery; it does not implement the theoretical bit-step
dovetail or claim its FPT runtime. Failure of these sufficient tests, or a resource limit, can leave the
result inconclusive, and general growth-only polynomial boundary
output remains unresolved.

\subsection{New bounded comparisons and certificate checks}
The second-stage main suite has 73 configurations. Every solve runs
sequentially in a fresh process with one numerical-library thread, a
two-second cooperative solver budget, a five-second hard worker deadline,
and a 512 MiB address-space cap. Replay runs in a separate five-second
worker. The table cap is 30,000 states per stage. Saved timings separate
solving, replay, preprocessing and total subprocess time; the latter
includes both workers. Source snapshots and final hashes identify the
implementation actually measured. Earlier provisional runs are archived
and excluded from the following totals.

\begin{center}
\begin{tabular}{lrrr}
\toprule
Lane & Configurations & Requests completed & Valid proof replays\\
\midrule
Archived grid baseline & 16 & 12 & 16\\
Completed grid & 16 & 12 & 16\\
Exact rational output & 9 & 9 & 9\\
Integrated recourse & 11 & 10 & 11\\
TU constraints & 16 & 13 & 15\\
Optimal-set descriptions & 5 & 5 & 5\\
\midrule
Total & 73 & 61 & 72\\
\bottomrule
\end{tabular}
\end{center}

Completion means the request made in that lane: the stated approximate
gap, an exact value, infeasibility, or a checked optimal-set description.
Partial bounds can pass replay without completing the request. The one
configuration without a proof is an intentionally invalid TU input,
correctly rejected before solving. Across the main suite, 54 enclosures
contain independent exact reference values and 29 configurations certify
exact values. The nine exact-output runs cover eight distinct models;
a separable case repeated without convex presolve already has an exact
initial bound, so that repeat is not evidence of nontrivial reconstruction.
The fresh random path, band and mixed cases provide such evidence.

Both grid versions finish new random paths through 64 variables and a
32-variable random tree. They finish the same 12 of 16 requests, with
summed solve-plus-replay subprocess times 15.18 and 15.34 seconds. This
comparison jointly changes minimum-degree to minimum-fill decomposition,
the message engine and sparse polishing; it establishes neither a speedup
nor an isolated cause for a timing difference. Minimum fill reduces the
supplied widths of the two full binary QPLIB instances from 25 to 19 and
95 to 93. Both still exceed the table budget, and neither public input
reaches its target. The three continuous bound-only QPLIB entries screened
from metadata have 14,400--39,204 variables; one has unbounded coordinates.
They exceed the prototype's declared 512-variable screen. These metadata
refusals are not solver runs or successful continuous public benchmarks.

The integrated recourse tests operate on the original models. On a
shuffled 17-variable affine star, automatic preprocessing removes 16
coordinates and certifies value $-1$ in 0.065 seconds of solve time and
0.047 seconds of replay. The 33-variable dense minimum-cut example reaches
gap $1/1024$ in 0.089 seconds, with 0.099 seconds for replay. These are
controlled structural examples. The two changing-response examples with
stiffness $M=1$ and $M=100$ each use 17 table states, six levels and seven
conditional QP queries, with certified gap $27/65536$. Their identical
state counts illustrate this certified curvature reduction; two parameter
values do not establish an empirical complexity law.

Two same-model constrained comparisons expose the value of the added
structure. For $x_0(1-x_0)+(x_1-1/3)^2+(x_2-1/3)^2$ with $x_1=x_2$,
coordinate hulls hit the per-stage cap after 15 levels and 33,502 total
processed states. Union filtering returns exact value zero after 43
levels and 2,538 total states, using two stationary-recovery LPs and
24 pivots; exhaustive face search is disabled. For
$1000(x_0-x_1)^2+(x_0-1/3)^2$ with $x_0=x_1$, the ambient-curvature
mode times out after 252,192 total states. The certified equality-tangent
bound yields exact value zero with 1,938 states. These small examples
exercise specific representation and curvature choices, rather than
establishing general solver superiority.

Eleven separately frozen extension configurations exercise polynomial,
implicit-boundary and mixed-submodular functionality. Seven reach their
requested outcome; three retain checked intervals after a table limit,
a cut limit or an unsupported signed class. One bounded symmetric-quartic
boundary search remains inconclusive and emits no exact-output certificate.
All ten emitted certificates replay. Eight numeric intervals enclose
independent exact references, and two implicit patches contain independently
known analytic optimizers.

The non-growth polynomial $(x-1/3)^4$ reaches its requested gap in
36 stages and 545 states. A mixed polynomial needs eight stages and
469 states. The native-integer example, with one fixed rational coordinate,
certifies value $-53/6$ after two stages and 17 states: its raw gap $1/24$
is strictly below its value-lattice spacing $1/18$. Independent enumeration
checks all 100 native assignments. The unplanted six-variable signed mixed
QP returns $-509/105$ using nine conditional QP queries, five greedy cuts
and 88 pivots. A separate fixture requires a genuine mixture of two bases;
a single-base lower bound $-13/16$ improves to exact value zero.

Across both new suites there are therefore \textbf{84 configurations,
82 valid proof replays and 68 completed requests}. Fourteen other runs
retain checked incomplete or unsupported-class bounds. The invalid TU
input and inconclusive boundary search are the two outcomes without
certificates. Reference evidence comprises 62 numeric enclosures and two
implicit-optimizer containment checks. The extension workers use the same
hard time and memory limits and total 1.584 seconds of subprocess time;
no process reaches a hard deadline or address-space cap. These figures
exclude provisional archives and the unchanged 74 historical runs.

The final combined topic-specific test command passes 158 tests in
4.114 seconds. Its saved output and 35-source hash manifest accompany the
completion package. These tests cover solver contracts, replay, invalid
inputs, resource exits and certificate corruption. They are separate from
the bounded benchmark configurations, and neither is a project-wide or
CI verification claim.
